% \iffalse meta-comment
% !TeX program  = XeLaTeX
% !TeX encoding = UTF-8
%
% Copyright (C) 2021 
% by Nanjing University Linux User Group <nju.lug@yaoge123.cn>
% 
% It may be distributed and/or modified under the conditions of the
% LaTeX Project Public License (LPPL), either version 1.3c of this
% license or (at your option) any later version.  The latest version
% of this license is in the file
%
%    https://www.latex-project.org/lppl.txt
%
% -----------------------------------------------------------------------
%
% The development version of the template can be found at
%
%    https://github.com/nju-lug/NJUThesis
%
% for those people who are interested.
%
%<*internal>
\iffalse
%</internal>
%
%<*readme-zh>
# NJU Thesis

[![njuthesis](https://img.shields.io/badge/njuthesis-latex-blue)](https://git.nju.edu.cn/nju-lug/nju-latex-templates)
[![overleaf](https://img.shields.io/badge/overleaf-supported-brightgreen)](https://tex.nju.edu.cn)
![build status](https://github.com/nju-lug/NJUThesis/actions/workflows/build.yml/badge.svg)

> 此项目是基于[NJU thesis 2021](https://github.com/FengChendian/NJUThesis2021)的重构版本，目前本科生版本**已经基本重构完成**，硕博论文**正在**根据研究生院的文件重构，但研究生院的规定并不完善，因此部分格式可能会参考本科的，**如果院系另有要求，请自行修改或者提交issue**。

原来的[NJU thesis 2021](https://github.com/FengChendian/NJUThesis2021)模板太过古老并且不支持硕士、博士论文，目前难以维护，因此为了模板的可持续发展，现在将对该项目进行重构，项目挂载于[NJU-LUG](https://github.com/nju-lug)组织仓库下，目前已[发布在CTAN](https://ctan.org/pkg/njuthesis)。

相关文档请参见[本项目wiki](https://github.com/nju-lug/NJUThesis/wiki)。

## 快速开始

从0.9.0版以后，本模板以DocStrip(`.dtx`)格式进行维护。常规的`.cls`格式文件可从Releases页面下载或使用以下方法提取。

### 关于DocStrip文件

DocStrip格式合并了LaTeX源代码和说明文档，更有利于宏包的发布。要使用该文件，请在当前工作目录下打开终端：

- 输入`xetex njuthesis.dtx`以提取宏包内容
- 输入`latexmk njuthesis.dtx`以生成并查看内置的文档

请注意，由于文件重名原因，项目中的示例论文模板已被重命名为`njuthesis-sample.tex`。

### 本地编译

1. [从南大镜像站下载](https://mirror.nju.edu.cn/download/app/TeX%20%E6%8E%92%E7%89%88%E7%B3%BB%E7%BB%9F)并安装TeXLive或者MiKTeX最新版
2. 下载包含[模板全部文件](https://github.com/nju-lug/NJUThesis/releases/latest)的压缩包`njuthesis.zip`
3. 根据需要更改`njuthesis-sample.tex`文件中的内容
4. 在模板根目录下运行`latexmk`编译文件，得到对应的PDF

#### **关于本地编译**

- 模板支持`latexmk`、`xelatex`、`lualatex`三种编译方式
- 如果使用`latexmk`，请务必使用完整的命令`latexmk -xelatex`以免产生不必要的问题
- 如果使用`xelatex`或`lualatex`，需要按照`xelatex/lualatex` -> `biber` -> `xelatex/lualatex` -> `xelatex/lualatex`的顺序编译四次
- 仓库中`.vscode`附带一份VScode的LaTeX WorkShop的配置文件，可根据需要使用

### [南大TeX](https://tex.nju.edu.cn)编译

1. 下载包含[模板全部文件](https://github.com/nju-lug/NJUThesis/releases/latest)的压缩包`njuthesis.zip`
2. 登录[南大TeX](https://tex.nju.edu.cn)，点击New Project -> Upload Project上传刚刚得到的压缩包，上传后`njuthesis-sample.tex`、`njuthesis.cls`等文件应在根目录，0.9.0以后的目录结构如下所示：

    ```shell
    NJUThesis-master/
    │
    ├─figure/
    │   njulogo.pdf
    │   njuname.pdf
    │
    │ njuthesis.bib
    │ njuthesis.cls
    │ njuthesis-sample.tex
    ```

3. 在南大TeX项目内页面左上角的`Menu`中，将编译器改为`XeLaTeX`
4. 编写TeX文档

### 其他

更多内容请参见项目Wiki中的[使用说明](https://github.com/nju-lug/NJUThesis/wiki/%E4%BD%BF%E7%94%A8%E8%AF%B4%E6%98%8E)。

你也可以下载[Release](https://github.com/nju-lug/NJUThesis/releases/latest)中附带的`njuthesis.pdf`，里面详尽阐述了模板的特性和使用方法。

## 常见问题

请参见[常见问题](https://github.com/nju-lug/NJUThesis/wiki/%E5%B8%B8%E8%A7%81%E9%97%AE%E9%A2%98)。

## 贡献

如果你发现了模板的问题，欢迎在Github中[提交issue](https://github.com/nju-lug/NJUThesis/issues)或者PR。

## 许可

本模板的发布遵守 LaTeX Project Public License（版本 1.3c 或更高）。

%</readme-zh>
%
%<*readme-en>
# The njuthesis class

## Overview

This package is intended for typesetting Nanjing University thesis with LaTeX, providing support for bachelor, master, and doctoral thesis. Compilation of this class requires either `xelatex` or `lualatex` engine.

## License

-----
    This work may be distributed and/or modified under the conditions of
    the [LaTeX Project Public License](http://www.latex-project.org/lppl.txt),
    either version 1.3c of this license or (at your option) any later
    version.
-----

## Contributing

If you run into any problem with the template, please contact us via Github issue or submit Pull Request.

## Wiki

Please see [Chinese wiki](https://github.com/nju-lug/NJUThesis/wiki)。

-----

Copyright © NJU-LUG
%</readme-en>
%
%<*internal>
\fi
\begingroup
  \def\NameOfLaTeXe{LaTeX2e}
\expandafter\endgroup\ifx\NameOfLaTeXe\fmtname\else
\csname fi\endcsname
%</internal>
%
%<*install>
\input l3docstrip.tex
\keepsilent
\askforoverwritefalse

\preamble

Copyright (C) 2021 
by Nanjing University Linux User Group <nju.lug@yaoge123.cn>

This file may be distributed and/or modified under the conditions of
the LaTeX Project Public License, either version 1.3 of this license
or (at your option) any later version.  The latest version of this
license is in:

   http://www.latex-project.org/lppl.txt

and version 1.3 or later is part of all distributions of LaTeX version
2005/12/01 or later.

To produce the documentation run the original source files ending with `.dtx'
through XeTeX.
    
\endpreamble

\generate{
  \usedir{tex/latex/njuthesis}
    \file{\jobname.cls}        {\from{\jobname.dtx}{class}}
%</install>
%<*internal>
  \usedir{source/latex/njuthesis}
    \file{\jobname.ins}        {\from{\jobname.dtx}{install}}
%</internal>
%<*install>
  \usedir{doc/latex/njuthesis}
  \nopreamble\nopostamble
    \file{README_EN.md}        {\from{\jobname.dtx}{readme-en}}
    \file{README.md}           {\from{\jobname.dtx}{readme-zh}}
}

\obeyspaces
\Msg{*************************************************************}
\Msg{*                                                           *}
\Msg{* To finish the installation you have to move the following *}
\Msg{* files into a directory searched by TeX:                   *}
\Msg{*                                                           *}
\Msg{* The recommended directory is TDS:tex/latex/njuthesis      *}
\Msg{*                                                           *}
\Msg{*     njuthesis.cls                                         *}
\Msg{*     njuthesis.ins                                         *}
\Msg{*     README.md                                             *}
\Msg{*     README_EN.md                                          *}
\Msg{*                                                           *}
\Msg{* To produce the documentation, run the file njuthesis.dtx  *}
\Msg{* through XeLaTeX.                                          *}
\Msg{*                                                           *}
\Msg{* Happy TeXing!                                             *}
\Msg{*                                                           *}
\Msg{*************************************************************}

\endbatchfile
%</install>
%
%<*internal>
\fi
%</internal>
%
%<class>\NeedsTeXFormat{LaTeX2e}
%<class>\RequirePackage{expl3}
%<!(readme-zh|readme-en)> \GetIdInfo  $Id: njuthesis.dtx 0.9.3 2021-09-19 20:00:00 +0800  NJU LUG <nju.lug@yaoge123.cn> $
%<class>  {Thesis template for Nanjing University}
%<class>\ProvidesExplClass{\ExplFileName}{\ExplFileDate}{\ExplFileVersion}{\ExplFileDescription}
%
%<*driver>
\ProvidesFile{njuthesis.dtx}
\documentclass{ctxdoc}
\usepackage{listings,xcolor,tabularray}
\UseTblrLibrary{booktabs,siunitx,diagbox}
\DefTblrTemplate{caption-tag}{default}{表\hspace{0.25em}\thetable}
\SetTblrStyle{caption-tag}{font=\bfseries}
\DefTblrTemplate{caption-sep}{default}{\quad}
\definecolor{njuviolet}{cmyk}{0.5,1,0,0.4}
\definecolor{njumagenta}{cmyk}{0.05,1,0.55,0}
\definecolor{njublue}{cmyk}{0.8,0.5,0,0}
\definecolor{njuyellow}{cmyk}{0,0.3,1,0}
\begin{document}
  \DocInput{njuthesis.dtx}
\end{document}
%</driver>
% \fi
%
% \title{\color{njuviolet}{The \textsc{NJuThesis} class\\ 南京大学学位论文模板}}
% 
% \author{^^A
% Nanjing University Linux User Group
% \thanks{E-mail: \href{mailto:nju.lug@yaoge123.cn}{nju.lug@yaoge123.cn}}}
%
% \date{v0.9.3 \\ Released 2021-09-19}
%
% \maketitle
%
% \begin{abstract}
% 南京大学学位论文\hologo{LaTeX}模板基于本科生院的论文撰写规范制作，同时参考研究生院提供的硕士、博士学位材料包，用于生成符合南京大学学位论文排版要求和相应的国家规范、行业标准的学位论文，旨在为同学提供毕业论文书写的方便。
% \end{abstract}
%
% \def\abstractname{Abstract}
% \begin{abstract}
% The \textsc{NJuThesis} class is intended for typesetting Nanjing University thesis with \hologo{LaTeX}, providing support for bachelor, master, and doctoral thesis. 
% \end{abstract}
%
% \vspace{1cm}
% \def\abstractname{特别声明}
% \begin{abstract}
% 请注意，本模板仍未完成开发，出现bug或者文档不完善属于正常情况，提issue或Pull Request即可。
% \href{https://git.nju.edu.cn/nju-lug/lug-introduction}{NJU LUG}始终欢迎您的加入与贡献！
% \end{abstract}
%
% \clearpage
%
% \setcounter{tocdepth}{4}
% \tableofcontents
% \clearpage
%
% \EnableDocumentation
%^^A \DisableDocumentation
% 
% \begin{documentation}
%
%
% \section{模板介绍}
% \textsc{NJuThesis}，即南京大学学位论文模板，是一个由南京大学LUG组织维护的，用于处理本校学生毕业论文排版需求的\hologo{LaTeX}模板。
%
% \subsection{历史沿革}
%
% 十几年来，多位热心校友发布过自己编写的模板，在GitHub上可考的包括
% \begin{enumerate}
%  \item 杨文博（\href{http://yangwenbo.com/}{@solrex}）的\href{https://github.com/solrex/njuthesis}{南京大学学位论文\hologo{LaTeX}模板}（2010） 
%  \item \href{https://github.com/fireblue}{@fireblue}基于solrex模板的的\href{https://github.com/fireblue/NJUThesis}{南京大学学位论文\hologo{LaTeX}模板}（2013）
%  \item \href{https://github.com/wenhai-zheng}{@wenhai-zheng}的\href{https://github.com/wenhai-zheng/NJUThesis}{NJUThesis}（2013）  
%  \item 曹增乐（\href{https://github.com/ZLCao}{@ZLCao}）的\href{https://github.com/ZLCao/NJUBachelor}{NJUBachelor}（2013-2016）
%  \item 胡海星（\href{http://haixing-hu.github.io/}{@Haixing-Hu}）的\href{https://github.com/Haixing-Hu/nju-thesis}{NJU-Thesis}（2013 - 2018）
%  \item 张楚珩（\href{https://github.com/zhangchuheng123}{@zhangchuheng123}）基于胡海星模板的\href{https://github.com/zhangchuheng123/NJUThesis}{NJUThesis}（2016）
%  \item 蒋炎岩（\href{http://ics.nju.edu.cn/~jyy/}{@jiangyy}）的\href{https://github.com/jiangyy/njuthesis}{南京大学山寨\hologo{LyX}研究生毕业论文模板}（2017）
%  \item \href{https://github.com/njuHan}{@njuHan}基于胡海星模板的\href{https://github.com/njuHan/njuthesis-nju-thesis-template/}{NJU-Thesis}（2018 - 2021）
%  \item 饶安逸（\href{https://anyirao.com/}{@AnyiRao}）基于张楚珩模板的\href{https://github.com/AnyiRao/NJUThesis2018/}{NJU Thesis 2018}（2018）
%  \item 赵懿晨（\href{https://fengchendian.github.io/about.html}{@FengChendian}）基于饶安逸模板的\href{https://github.com/FengChendian/NJUThesis2021}{NJU Thesis 2021}（2021）
% \end{enumerate}
% 至于不幸而未进行代码版本管理的，甚至于说以压缩包形式流传于各人硬盘中的，就更不可计数了。
%
% 与其让一千个学生做一千种不同的模板，不如有人牵起头来，集中力量办大事。因此，南京大学 Linux User Group 的有志之士在2021年暑期联合起来，参考以上的现成模板，兼顾友校thuthesis、fduthesis等优秀项目，构建了这一全新模板，力求通过\hologo{LaTeX3}语法和清淅的接口实现良好的易用性和可维护性。
%
% 目前，本模板的主要维护者包括：
%
% \begin{itemize}
%     \item 赵懿晨
%     \item 熊煜
% \end{itemize}
%
% \subsection{君\hologo{LaTeX}本当上手}
% \hologo{LaTeX}并不是一种零门槛的语言，亦不是南京大学学生的必备技能。有同学听到“\emph{给毕业论文交差也得写代码}”这种事情就头疼，实在是我们不愿意见到而必须要面对的情况。目前，中文互联网上的\hologo{LaTeX}写作教程不能说屈指可数，起码也是汗牛充栋，比较著名的包括刘海洋的《\hologo{LaTeX}入门》等等都是不错的参考资料。然而如何快速准确找到所需要的用法，尤其是\hologo{LaTeX3}的新接口，需要一定的时间或者良好的英文阅读能力。在本文档的写作中，我们无力从盘古开天辟地开始从头构建您的相关知识体系，只能\textbf{尽力做到有求必应}，在需要的地方顺带提一嘴，保证您少走点弯路。
%
% \subsection{关于本说明手册}
%
% 本模板通过\textsc{DocStrip}进行维护，实际上是一种\emph{文学编程}。文学编程的本体是那些用来说明的文字，按写书写文章的方式组织。
% \file{njuthesis.dtx}包含njuthesis模板的所有信息。前半部分为模板说明，即各种常用命令与必要的使用方法；后半部分为代码实现，为带有详细注释的\hologo{TeX}源代码。后者配合附录，为有意了解本模板细节的同学提供简明的指引。
%
% 下文中将出现若干特殊格式，用来指示代码组件。譬如，\file{file-example.xxx}为文件名，\pkg{pkg-example}为宏包名，\opt{opt-example}为选项名，\env{env-example}为需要|\begin{}|和|\end{}|的环境名，等。
%
% 如果您在使用最新版模板时发现了任何问题，抑或有意向参与本模板的维护，请通过GitHub issue或者QQ群聊联系LUG。
%
% \section{安装}
% \label{sec:setup}
% 本节主要介绍本宏包的获取、安装以及编译方式
% 
% \subsection{下载模板}
%
% 本宏包已经\href{https://ctan.org/pkg/njuthesis}{发布在CTAN}，将在明年进入\hologo{TeX}\,Live 2022 和\hologo{MiKTeX}发行版。然而CTAN版更新不及时，更推荐使用以下方式下载。
%
% 
% \subsubsection{发布版}
%
% 从本模板的\href{https://github.com/nju-lug/NJUThesis/releases/latest}{Github Releases页面}下载最新的发布版。
% 如果Github访问速度较慢，也可以使用\href{https://mirror.nju.edu.cn/github-release/nju-lug/NJUThesis/LatestRelease/}{南大镜像}
%
% \subsubsection{源代码}
%
% 从本模板的\href{https://github.com/nju-lug/NJUThesis/}{ Github 仓库}或者\href{https://git.nju.edu.cn/nju-lug/nju-latex-templates/njuthesis}{NJU Git 同步镜像}下载最新的dtx文件自行安装
%
% \subsection{本地编译}
%
% \subsubsection{安装\hologo{TeX}发行版}
%
% 首先需要下载\hologo{TeX}软件发行版，校园网环境中使用\href{https://mirror.nju.edu.cn/download/app/TeX%20%E6%8E%92%E7%89%88%E7%B3%BB%E7%BB%9F}{南大镜像站}可以获得最好的体验。\textbf{推荐使用最新的\hologo{TeX}\,Live 2021或者\hologo{MiKTeX} 21以避免潜在的兼容性问题。}
%
% \begin{itemize}
%     \item 为了避免不必要的麻烦，请尽可能下载 full 版本，如 texlive-full。简而言之，下载大的那个。
%     \item 并且，尽可能使用最新版（截至目前是 2021）。2020 及之前版本使用 PDF 格式的图片可能会出现加粗问题以及l3语法解释错误。
% \end{itemize}
%
% 下表是目前经过测试的环境。如果有其他可用不可用的环境，欢迎补充。
% \begin{table}[ht]
%     \centering
%     \caption{经过测试的环境}
%     \label{tab:1}
%     \begin{tabular}{ccc}
%           \toprule
%           OS & TeX & 测试情况 \\
%           \midrule
%           Windows 10   & \hologo{TeX}\,Live 2020 & cref存在格式问题  \\
%           Windows 10   & \hologo{TeX}\,Live 2021 & 通过 \\
%           Windows 10   & \hologo{MiKTeX}         & 通过 \\
%           macOS 10.15  & mac\hologo{TeX} 2021    & 通过 \\
%           Ubuntu 20.04 & \hologo{TeX}\,Live 2021 & 通过 \\
%           Arch Linux   & \hologo{TeX}\,Live 2021 & 通过 \\
%           \bottomrule
%     \end{tabular}
% \end{table}
%
% \subsubsection{选择编辑器}
%
% 配置完编译器后，还需要一个\textbf{文本编辑器}作为前端来完成\texttt{.tex}文件内容的写作。
%
% 至今仍有相当一部分人认为Windows自带的\textit{记事本}是最好的文本编辑器，但对于本项目而言，在此诚心诚意地推荐你使用\textbf{更现代更美观更多功能}的编辑器，譬如\emph{安装了 LaTeX Workshop 插件 的 \href{https://code.visualstudio.com/}{Visual Studio Code}}，来完成论文编写。你也可以根据个人的喜好随便使用其他编辑器，如 TeXworks、TeX Studio 等，顺手就行。
%
% 若使用 LaTeX Workshop 插件，本项目在|.vscode/|中提供一份简易配置，可以省略初始配置步骤直接使用。
%
% \subsubsection{编译顺序}
% 假设即将使用的文件名为\file{mynjuthesis.tex}
% 应采用以下命令顺序进行编译，以生成正确的目录、编号和参考文献条目。
% \begin{ctexexam}
%   xelatex mynjuthesis
%   biber mynjuthesis
%   xelatex mynjuthesis
%   xelatex mynjuthesis
% \end{ctexexam}
% 使用 \pkg{latexmk}
% \begin{ctexexam}
%   latexmk -xelatex mynjuthesis
% \end{ctexexam}
%
% \begin{ctexexam}
%   lualatex mynjuthesis
%   biber mynjuthesis
%   lualatex mynjuthesis
%   lualatex mynjuthesis
% \end{ctexexam}
% 使用 \pkg{latexmk}
% \begin{ctexexam}
%   latexmk -lualatex mynjuthesis
% \end{ctexexam}
%
% 编译产物为\file{njuthesis.pdf}，位于主目录下。此外还会生成一系列中间文件，可以选择使用以下命令进行清理。
%
% \begin{ctexexam}
%   latexmk -c
% \end{ctexexam}
%
% \subsection{在线编译}
%
% 相信你在接触了本地编译以后，很快就会意识到一些十分显然的事实，譬如\hologo{TeX}编译器安装过程较为漫长，占用空间过大，而且在一部分处理器性能不佳的电脑上需要较长编译时间\footnote{其实这三点都是对广大的Windows用户说的，同一个模板在Linux编译可以节省一半耗时}。拒绝接受这些麻烦的同学不妨尝试本节介绍的在线编译方法。
%
% \subsubsection{南大\hologo{TeX}平台简介}
%
% \href{https://tex.nju.edu.cn}{南大\hologo{TeX}}基于开源的ShareLaTeX平台\footnote{理论上在\href{https://doc.nju.edu.cn/books/latex}{这个网站}能找到一段平台简介，实际上大家都有意无意地鸽了，下次一定补上。}，于2021年3月4日正式上线，面向南京大学全体师生开放，首次使用需凭学校邮箱自助注册账号。
%
% \subsubsection{操作步骤}
%
% \begin{enumerate}
%     \item 下载模板全部文件
%     \item 访问\href{https://tex.nju.edu.cn}{南大\hologo{TeX}}，点击界面右上方Register，使用\emph{南京大学邮箱}注册账号并登录
%     \item 点击New Project -> Upload Project上传刚刚得到的zip文件，上传后njuthesis.tex、njuthesis.cls等文件应在根目录，目录结构如{{sec:directory}}所示
%     \item 在项目页面左上角的Menu中，将编译器改为\hologo{XeLaTeX}或者\hologo{LuaLaTeX}
%     \item 编写论文
%     \item 点击Compile按钮进行编译和预览
%     \item 点击编译按钮右侧第三个按钮下载产物
% \end{enumerate}
%
% \subsubsection{关于Overleaf平台}
%
% 由于\href{https://www.overleaf.com/}{Overleaf平台}的\hologo{TeX}\,Live版本停留在2020，\texttt{cleveref}包在引用章节时会生成错误的标签，引发格式错误；而南大\hologo{TeX}通过及时更新规避了这一问题。因此\emph{请务必不要使用Overleaf官网进行编译}。
%
%
%
% \section{使用方法}
%
% \subsection{模板选项}
%
% 尖括号内列出了若干个允许的选项，其中加粗的为默认选项。
%
% 封面标题留空即可自动化换行
%
% 第二导师不填则留空，目前来说拥有两位指导老师并不是一种普遍状况。
%
% \begin{function}[added=2021-09-07]{nlcover}
%   \begin{syntax}
%     nlcover = <\TTF>
%   \end{syntax}
%
% 通过此项可选择是否添加国家图书馆封面，本选项仅对于研究生有效
% \end{function}
%
% 
% \subsubsection{学位信息}
% \begin{function}[added=2021-09-07]{degree}
%   \begin{syntax}
%     degree = <(ug)|mg|mf|phd>
%   \end{syntax}
%
% 选择学位，可选：
% \opt{ug}（默认本科），\opt{mg}即学术硕士，\opt{mf}，\opt{phd}。
% \end{function}
%   
%
% \begin{function}[added=2021-09-07]{type}
%   \begin{syntax}
%     type = <(thesis)|design>
%   \end{syntax}
%
% 用于选择文档类型并将相关字段输出在封面和摘要页。可选值包括\opt{thesis}（表示毕业论文）和\opt{design}（表示毕业设计），缺省值为thesis。
% 
% \end{function}
%   
%
%
% \subsubsection{指定字体}
%
% 学校论文格式要求使用的字体一般已经预装在各个操作系统，本模板针对不同平台进行了自动检测适配，可以开箱即用。
%
% 如果希望更改本模板使用的字体，请填写以下两个选项以覆盖操作系统检测得到的设置。
% 
% \begin{function}[added=2021-09-07]{customchinesefont}
%   \begin{syntax}
%     customchinesefont = <(windows)|macos|fandol|founder|noto|none>
%   \end{syntax}
%
% 手动指定中文字体，可选：
% \opt{windows}（默认）， \opt{macos}， \opt{fandol}， \opt{founder}，\opt{noto}，\opt{none}。
% \end{function}
%
% \begin{function}[added=2021-09-07]{customlatinfont}
%   \begin{syntax}
%     customlatinfont = <(windows)|macos|gyre|none>
%   \end{syntax}
%
% 手动指定西文字体，可选：
% \opt{windows}（默认）， \opt{macos}， \opt{gyre}， \opt{none}。
% \end{function}
%
% 根据学校论文格式的要求，本模板使用的中文字体主要有{\songti 宋体}、{\heiti 黑体}、{\kaishu 楷体}、{\fangsong 仿宋}四种；
% 西文字体主要有衬线体（\textup{serif}）、无衬线体（\textsf{sans-serif}）和等宽字体（\texttt{monospace}）三种。
% 在生成论文时，需要将每一种字体（宋体、楷体）与一种字库（操作系统中的 ttf 文件）相对应。
% 由于不同操作系统自带的字库不同，在不同的操作系统上使用同一套配置会导致缺字、编译出错等后果。
% 因此，本模板对于中文字体和西文字体都提供了多种字体配置。
%
% 本模板提供的中文字体配置如表 \ref{tab:chinesefontset} 所示。
% 在不指定字体配置的情况下，本模板默认使用与操作系统相对应的字体配置（见表中前三行，Windows 和 macOS 以外的系统采用 Fandol 配置）。
% 此外，我们也单独提供了方正和思源（\href{https://www.google.com/get/noto/#serif-hans}{思源宋体}、\href{https://www.google.com/get/noto/#sans-hans}{思源黑体}）两套中文字体配置。
%
% \begin{table}[htbp]
%   \centering
%   \begin{talltblr}[
%    caption = {中文字体配置列表},
%    label = {tab:chinesefontset},
%    note{a} = {使用方正小标宋作为方正书宋的粗体，使用方正黑体作为方正细黑一的粗体},
%    note{b} = {思源字体并不包含楷书和仿宋，而Adobe楷体和仿宋难以直接下载，因此使用方正字体代替；考虑到获取字体的难易程度，采用 \href{https://www.google.com/get/noto}{noto} 字体表示思源字体},
%   ] { cell{2}{2,4} = {njuyellow}, cell{4}{4} = {njuyellow}, colspec = {ccccc} }
%     \toprule
%     配置名称 & 宋体 & 黑体 & 楷书 & 仿宋 \\
%     \midrule
%     Windows          & 中易宋体    & 中易黑体    & 中易楷体    & 中易仿宋 \\
%     macOS            & 华文宋体    & 华文黑体    & 华文楷体    & 华文仿宋 \\
%     \href{https://www.ctan.org/pkg/fandol} {Fandol}           
%                      & Fandol宋体  & Fandol黑体  & Fandol楷体  & Fandol仿宋 \\
%     \href{https://www.foundertype.com}     {方正}\TblrNote{a} 
%                      & 方正书宋    & 方正细黑一  & 方正楷体    & 方正仿宋 \\
%     \href{https://github.com/adobe-fonts}  {思源}\TblrNote{b} 
%                      & 思源宋体    & 思源黑体    & 方正楷体    & 方正仿宋 \\
%     \bottomrule
%   \end{talltblr}
% \end{table}
%
% 由于各操作系统对自带字体的支持程度不同，本模板对不同字体加粗的方式有所差别，表 \ref{tab:chinesefontset} 中，黄色底色的字体采用算法加粗。
% 除此以外，各种字体在设计上也存在差异，因此本模板在不同操作系统下编译的结果可能存在细微不同，使用时请注意。
% 若想要获得与 Word 版本模板最为贴近的体验，可以通过手动安装 Windows 自带的中易宋体、中易黑体、中易楷体和中易仿宋，并设置使用 Windows 字体来实现。
%
% 本模板提供的西文字体配置如表 \ref{tab:latinfontset} 所示。
% 在不指定字体配置的情况下，本模板默认使用与操作系统相对应的字体配置（见表中前三行，Windows 和 macOS 以外的系统采用 \opt{gyre} 配置）。
%
% \begin{table}[htbp]
%   \centering
%   \caption{西文字体清单}
%   \label{tab:latinfontset}
%   \begin{tabular}{cccc}
%       \toprule
%       配置名称 & 衬线体 & 无衬线体 & 等宽字体 \\
%       \midrule
%       Windows  & Times~New~Roman           & Arial                   & Courier~New \\
%       macOS    & Times~New~Roman           & Arial                   & Menlo \\
%       \href{http://www.gust.org.pl/projects/e-foundry/tex-gyre}{gyre} 
%                & \Hologo{TeX}~Gyre~Termes  & \Hologo{TeX}~Gyre~Heros & \Hologo{TeX}~Gyre~Cursor \\
%       \bottomrule
%   \end{tabular}
% \end{table}
%
% 若要手动选择字体，则需将字体设置为 \opt{none}，再通过 \pkg{fontspec} 宏包或 \pkg{ctex} 宏包中给出的命令来实现。
%
% 最后，出于统一格式的需要，本模板还选取了与正文字体相似度最高的数学字体（\href{https://www.stixfonts.org}{STIX 字体}）用于公式排版。
%
% 本节提及的字库中，中文字库只有方正书宋、方正黑体、方正楷体、方正仿宋、两种思源字体和四种Fandol字体能够\emph{免费商用}；
% 英文字库中只有 Times New Roman、STIX 和三种\Hologo{TeX} Gyre字体能够\emph{免费商用}。
% 因此，若需要将本模板用于商业相关的活动，请咨询专业人士。
%
% \subsection{个人信息}
%
% \begin{function}{\njusetup}
% 
% 定义了用于文档封面的诸多属性参数，
% 写作时修改相应字符串即可。
%
% \end{function}
%
% 本科生可以参考如下列表输入个人信息。注意不能有空行，有空格的内容需要将空格替换为波浪线或者两端加上大括号
%
% \begin{ctexexam}
%   \njusetup {
%       info = {
%           TitleA = 第一行标题,
%           TitleB = 第二行标题,
%           Title* = {My title in English}, 
%           Grade = 2018,
%           StudentID = 189114514,
%           StudentName = 周煜华,
%           StudentName* = Zhou~Yuhua,
%           Department = 拉太赫科学与技术学院,
%           Department* = School~of~\hologo{LaTeX},
%           Major = 拉太赫语言学,
%           Major* = \hologo{LaTeX}~Linguistics,
%           SupervisorA = 李成殿,
%           SupervisorA*= Li~Chengdian,
%           SupervisorATitle = 教授,
%           SupervisorATitle* = Professor,
%           SupervisorB = 孙赫弥,
%           SupervisorB* = Sun~Hemi,
%           SupervisorBTitle = 副教授,
%           SupervisorBTitle* = Associate professor,
%           SubmitDate = 2021年8月10日,
%       }
%   }
% \end{ctexexam}
%
% 研究生可以参考如下列表输入个人信息
%
% \begin{ctexexam}
%   \njusetup {
%       info = {
%           TitleA = 第一行标题,
%           TitleB = 第二行标题,
%           TitleC = 第三行标题,
%           Title* = {My title in English}, 
%           Grade = 2018,
%           StudentID = dz18114514,
%           StudentName = 周煜华,
%           StudentName* = Zhou~Yuhua,
%           Department = 拉太赫科学与技术学院,
%           Department* = School~of~\hologo{LaTeX},
%           Major = 拉太赫语言学,
%           Major* = \hologo{LaTeX}~Linguistics,
%           Field = 拉太赫语言在现当代的使用, 
%           Field* = Contemporary~Usage~of~the~\hologo{LaTeX}~Language, 
%           SupervisorA = 李成殿,
%           SupervisorA*= Li~Chengdian,
%           SupervisorATitle = 教授,
%           SupervisorATitle* = Professor,
%           SubmitDate = 2021年8月10日,
%           SubmitDate* = {Aug 10, 2021}, 
%           DefendDate = 2021年9月19日,
%           ReviewerChairman = 张晓山~教授,
%           ReviewerA = 王瑞希~教授,
%           ReviewerB = 郭德纲~副教授,
%           ReviewerC = 华芈库~教授,
%           ReviewerD = 戴菲菲~教授,
%           Classification = 0175.2,
%           SecurityLevel = 限制,
%           UDC = 004.72,
%           SupervisorContact = {拉太赫科学与技术学院 枝江市结丘路 19 号 114514},
%       }
%   }
% \end{ctexexam}
%
%
%
% \subsection{编写正文}
%
% 与导言区相对的是正文，位于\env{document}环境内。
% 
% \begin{ctexexam}
%   \begin{document}}
%     text
%   \end{document}
% \end{ctexexam}
%
% 加上空行就可以另起一段。
%
% 
% 在页面布局方面，本科生无页眉，页面编号居中位于页脚；研究生无页脚，页眉包括章节名和页面编号（尚未确定）。
%
% \subsubsection{文字}
%
% 为了避免原生|\uline|命令在中文环境下不能正常换行的问题，本模板使用|\njuline|作为替代方案，分别对\hologo{XeLaTeX} 的\pkg{xeCJKfntef}包和\hologo{LuaLaTeX}的\pkg{lua-ul}包中的下划线命令进行了包装。
%
% \subsubsection{图片}
%
% 所有图片默认存放在主目录下的\texttt{figure/}文件夹内，允许的文件名包括\file{*.jpg}, \file{*.png}, \file{*.pdf}, \file{*.eps}, 
%
% 下面这个例子展示了如何插入一张名称为\file{myfig.jpg}，宽度占页面文字宽度一半的图片。注意标签命令|\label|必须放在图注命令|\caption|之后。
%
% \begin{ctexexam}
%   \begin{figure}
%       \includegraphics[width=0.5\textwidth]{myfig}
%       \caption{我的图}
%       \label{fig:myfig}
%   \end{figure}
% \end{ctexexam}
%
% 
% 文字环绕图像，\env{wrapfigure}后面不能有空行
%
% \begin{ctexexam}
%   \begin{wrapfigure}{r}{0cm}
%       \includegraphics[width=0.15\textwidth]{myfig}
%       \caption{我的图}
%       \label{fig:myfig}
%   \end{wrapfigure}
%   <text>
% \end{ctexexam}
% 
% TODO：subfigure
%
% \subsubsection{表格}
%
% 下面这个例子展示了如何插入一张表格。通过\pkg{booktabs}提供的|\toprule|、|\midrule|和|\bottomrule|，我们可以很轻松地绘制出一张漂亮规范的三线表。
% \begin{ctexexam}
%   \begin{table}[htbp]
%       \caption{我的表}
%       \label{tab:testtab}
%       \begin{tabular}{ccc}
%           \toprule
%           OS & TeX & 测试情况 \\
%           \midrule
%           南大TeX & Overleaf & √ \\
%           \bottomrule
%       \end{tabular}
%   \end{table}
% \end{ctexexam}
%
% TODO：\pkg{multirow}\pkg{multicol}合并单元格
%
% \subsubsection{代码}
%
% 由\pkg{listings}提供代码排版。代码块\env{lstlisting}，注意这个环境中的缩进空格会如实输出
% \begin{ctexexam}
%   \begin{lstlisting}
%   <code>
%   \end{lstlisting}
% \end{ctexexam}
%
% 行间代码|\lstinline!<code>!|，其中包裹代码的两个感叹号并不是绝对的，可以替换成任何两个相同的没有在这段代码中出现过的符号。
%
%
% \subsubsection{数学}
%
% 行内公式形如|$\mathrm{e}^{(a+2b)x}$|
%
% 带有编号的行间公式\env{equation}
% \begin{ctexexam}
%   \begin{equation}\label{eq:myeqlabel}
%       \pi
%   \end{equation}
% \end{ctexexam}
%
% 如不需要编号，可以用\env{equation*}或者 |\[<eq>\]|
%
% 本模板也提供了一系列的数学环境。证明环境会在结尾添加证毕符号
% \begin{ctexexam}
%   \begin{proof}
%       证明我是我
%   \end{proof}
% \end{ctexexam}
%
% 普通环境
% \begin{ctexexam}
%   \begin{definition}[他人]
%       定义他人即地狱
%   \end{definition}
% \end{ctexexam}
% 
% 
% \subsubsection{脚注}
% |\footnote{<text>}|
% 
% 
% \subsubsection{引用}
% 
% 使用\pkg{cleveref}宏包实现了带图、表等项目名称的智能引用，仅需使用|\cref{<label>}|命令。
%
% 超链接|\href{<url>}{<name>}|
% 
% 
% \subsection{特殊页面}
% 
% 本模板还提供一系列环境用于生成所需的特殊页面
%
% \subsubsection{封面}
%
% 封面可使用如下命令生成。
% 
% \begin{ctexexam}
%   \maketitle
% \end{ctexexam}
%
% 本科生仅会生成中文封面；研究生会生成中英文封面。如果选择了\opt{nlcover}，也会生成研究生的国家图书馆封面。
%
% \subsubsection{摘要页}
%
% 中文关键词，每个关键词之间用“；”分开，最后一个关键词不打标点符号
% 英文关键词。关键词之间用英文半角逗号隔开，末尾无符号。
% 
% \begin{ctexexam}
%   \begin{abstract}
%       我的中文摘要
%       \keywords{我；就是；充数的；关键词}
%   \end{abstract}
%
%   \begin{englishabstract}
%       My abstract in English
%       \englishkeywords{Dummy; Keywords; Here}
%   \end{englishabstract}
% \end{ctexexam}
%
% 在使用\hologo{LuaLaTeX}编译时，研究生中文摘要页的标题会出现空格无下划线的问题，目前正在积极寻求解决方法。
%
% \subsubsection{前言页}
%        
% 使用\env{preface}环境定义
% 
% \begin{ctexexam}
%   \begin{preface}
%       我的前言
%       \vspace{1cm}
%       \begin{flushright}
%       我的名字\\
%       时间地点
%       \end{flushright}
%   \end{preface}
% \end{ctexexam}
%
% \subsubsection{目录页}
%
% 目录、图片清单和表格清单可使用如下命令生成。
% 
% \begin{ctexexam}
%   \tableofcontents
%   \listoffigures
%   \listoftables
% \end{ctexexam}
%
% \subsubsection{致谢页}
%
% 同前言，使用\env{acknowledgement}环境
% 
% \begin{ctexexam}
%   \begin{acknowledgement}
%       感谢NJU LUG
%   \end{acknowledgement}
% \end{ctexexam}
%
% \subsubsection{参考文献页}
%
% 使用以下命令，\opt{bibintoc}可以将参考文献页插入目录
% 
% \begin{ctexexam}
%   \printbibliography[heading=bibintoc,title=参考文献]
% \end{ctexexam}
% 
% \subsubsection{附录页}
%
% 附录放在|\appendix|命令后，以英文字母进行编号，正常通过|\chapter{}|编写
%
% 是否需要索引？
% 
%
% \subsection{参考文献}
%
% 符合GB7714-2015规范。使用\hologo{biber}作为\hologo{BibTeX}后端。需要使用|biber|命令手动编译才会显示
%
% 默认参考文献存储于主目录下的\file{njuthesis.bib}，直接向其中粘贴新的参考文献即可。
%
% 如果希望额外添加参考文献列表，可以在主文件\file{mynjuthesis.tex}的导言区中多次调用|\addbibresource|命令：
%
% \begin{ctexexam}
%   \addbibresource{file1.bib}
%   \addbibresource{file2.bib}
% \end{ctexexam}
% 注意本命令与|\bibliography{bibfile1,bibfile2}|不同，不可以用逗号分隔多个输入文件，且必须使用带扩展名的完整文件名。
%
% \subsubsection{使用EndNote}
%
% 南京大学信息化建设管理服务中心已购买\href{https://itsc.nju.edu.cn/EndNote/list.htm}{EndNote}供全校师生免费试用，最新版为EndNote 20。
%
%
%
% 
%
% \subsubsection{使用Zotero}
%
% \href{https://www.zotero.org/}{Zotero}是一款免费的文献管理软件，支持所有桌面平台。
%
% 在保持Zotero程序运行的情况下，点击浏览器工具栏的Zotero Connector插件即可自动从网页抓取参考文件信息。Zotero也有知网插件
%
% \subsubsection{文段内引用}
%
% |\cite{myref}|，|\citeauthor{myref}|
%
% 
% \end{documentation}
%
% \begin{implementation}
%
% \section{代码实现}
% @@在l3docstrip中表示名空间，在生成cls时会被相应字段替换，譬如在\pkg{njuthesis}中|@@=nju|。
% 尖括号包裹的|<*class>||</class>|用来指定某段代码属于哪个文件。
%
%    \begin{macrocode}
%<@@=nju>
%<*class>
\LoadClass[
  a4paper,
  twoside,
  UTF8,
  scheme=chinese,
%    \end{macrocode}
%
% 关于行距，\hologo{LaTeX}默认1.2行距，word默认行距是1.3，要求1.5倍word行距，故
% \[ 1.5\times\frac{1.3}{1.2} = 1.625\]
%    \begin{macrocode}
  linespread=1.625,
%    \end{macrocode}
%
% 默认不载入任何字体，供模板自行设置
%    \begin{macrocode}
  fontset=none,
  zihao=-4
  ]{ctexbook}[2018/04/01]
%    \end{macrocode}
%
% 宏包\pkg{l3keys2e}用于配置模板选项。
%    \begin{macrocode}
\RequirePackage{l3keys2e}
%    \end{macrocode}
%
% 是否需要国家图书馆封面的设置
%    \begin{macrocode}
\keys_define:nn { nju }
{
  nlcover           .bool_set:N   =   \g_@@_nlcover_bool,
  nlcover           .initial:n    =   false,
}
%    \end{macrocode}
% 
%
% \begin{variable}{\l_@@_info_degree_tl,\l_@@_info_type_tl}
% 用于存储学位名称的变量，注意宏的命名，l代表局部变量，g代表全局变量
%    \begin{macrocode}
\tl_new:N \l_@@_info_degree_tl
\tl_new:N \l_@@_info_type_tl
%    \end{macrocode}
% \end{variable}
%
% 学位信息的设置
%    \begin{macrocode}
\keys_define:nn { nju }
{
  degree            .choices:nn   =
  { ug, mg, mg, phd }
  { \tl_set_eq:NN \l_@@_info_degree_tl \l_keys_choice_tl },  
  degree            .initial:n    =   ug,
  type              .choices:nn   =
  { thesis, design }
  { \tl_set_eq:NN \l_@@_info_type_tl   \l_keys_choice_tl }, 
  type              .initial:n    =   thesis,
}
%    \end{macrocode}
% 
% \subsection{个人信息}
% 输入个人信息的区域。
%    \begin{macrocode}
\keys_define:nn { nju }
{
  info.meta:nn = { nju / info } { #1 }
}
%    \end{macrocode}
%
% 题目
%    \begin{macrocode}
\keys_define:nn { nju / info }
{
  TitleA            .tl_set:N     =   \l_@@_info_title_a_tl,
  TitleB            .tl_set:N     =   \l_@@_info_title_b_tl,
  TitleC            .tl_set:N     =   \l_@@_info_title_c_tl,
  Title*            .tl_set:N     =   \l_@@_info_title_en_tl,
}
%    \end{macrocode}
%
% 年级、学号、姓名
%    \begin{macrocode}
\keys_define:nn { nju / info }
{
  Grade             .tl_set:N     =   \l_@@_info_grade_tl,
  StudentID         .tl_set:N     =   \l_@@_info_id_tl,
  StudentName       .tl_set:N     =   \l_@@_info_author_tl,
  StudentName*      .tl_set:N     =   \l_@@_info_author_en_tl,
}
%    \end{macrocode}
%
% 院系、专业、方向
%    \begin{macrocode}
\keys_define:nn { nju / info }
{
  Department        .tl_set:N     =   \l_@@_info_dept_tl,
  Department*       .tl_set:N     =   \l_@@_info_dept_en_tl,
  Major             .tl_set:N     =   \l_@@_major_tl,
  Major*            .tl_set:N     =   \l_@@_major_en_tl,
  Field             .tl_set:N     =   \l_@@_field_tl,
  Field*            .tl_set:N     =   \l_@@_field_en_tl,
}
%    \end{macrocode}
%  
% 导师
% 
%    \begin{macrocode}
\keys_define:nn { nju / info }
{
  SupervisorA       .tl_set:N     =   \l_@@_info_supv_a_tl,
  SupervisorA*      .tl_set:N     =   \l_@@_info_supv_a_en_tl,
  SupervisorATitle  .tl_set:N     =   \l_@@_info_supv_a_title_tl,
  SupervisorATitle* .tl_set:N     =   \l_@@_info_supv_a_title_en_tl,
}
%    \end{macrocode}
%
% 第二导师
%
%    \begin{macrocode}
\keys_define:nn { nju / info }
{
  SupervisorB       .tl_set:N     =   \l_@@_info_supv_b_tl,
  SupervisorB*      .tl_set:N     =   \l_@@_info_supv_b_en_tl,
  SupervisorBTitle  .tl_set:N     =   \l_@@_info_supv_b_title_tl,
  SupervisorBTitle* .tl_set:N     =   \l_@@_info_supv_b_title_en_tl,
}
%    \end{macrocode}
%
% 提交日期
%    \begin{macrocode}
\keys_define:nn { nju / info }
{
  SubmitDate        .tl_set:N     =   \l_@@_submit_date_tl,
  SubmitDate*       .tl_set:N     =   \l_@@_submit_date_en_tl,
}
%    \end{macrocode}
%
% 答辩
%    \begin{macrocode}
\keys_define:nn { nju / info }
{
  DefendDate        .tl_set:N     =   \l_@@_defend_date_tl,
  ReviewerChairman  .tl_set:N     =   \l_@@_info_chairman_tl,
  ReviewerA         .tl_set:N     =   \l_@@_info_reviewer_a_tl,
  ReviewerB         .tl_set:N     =   \l_@@_info_reviewer_b_tl,
  ReviewerC         .tl_set:N     =   \l_@@_info_reviewer_c_tl,
  ReviewerD         .tl_set:N     =   \l_@@_info_reviewer_d_tl,
}
%    \end{macrocode}
%
% 国家图书馆封面相关
%    \begin{macrocode}
\keys_define:nn { nju / info }
{
  Classification    .tl_set:N     =   \l_@@_info_classif_tl,
  SecurityLevel     .tl_set:N     =   \l_@@_info_seclv_tl,
  UDC               .tl_set:N     =   \l_@@_info_udc_tl,
  SupervisorContact .tl_set:N     =   \l_@@_info_supv_cont_tl,
}
%    \end{macrocode}
% 定义用于设置个人信息的命令
%    \begin{macrocode}
\NewDocumentCommand \njusetup { m }
{ \keys_set:nn { nju } { #1 } }
%    \end{macrocode}
%
% 使用\pkg{geometry}设置页边距。
%    \begin{macrocode}
\RequirePackage[    
    top=2.5cm,
    bottom=2.5cm,
    left=3.2cm,
    right=3.2cm
]{geometry}
%    \end{macrocode}
%
% 
%    \begin{macrocode}
\RequirePackage{xparse}
\RequirePackage[hyphens]{url} % generate better linebreaks in the url
\RequirePackage{dashundergaps}
\RequirePackage{setspace}
\RequirePackage{lastpage}
\RequirePackage{emptypage} % 清除空白页的页码
%    \end{macrocode}
%
% 
%    \begin{macrocode}
\RequirePackage{listings} % 代码环境
\RequirePackage{enumitem} % 用于修改列表环境
\RequirePackage{caption}
\RequirePackage{floatrow} % 用于图表等页面元素的定位
\RequirePackage{booktabs} % 用于绘制三线表
\RequirePackage{multirow} % Cells occupying multiple rows in tables
\RequirePackage{multicol} % Multiple columns in dictionary
\RequirePackage{siunitx} % 用于书写单位符号
%    \end{macrocode}
%
% 用于化学
%    \begin{macrocode}
\RequirePackage[version=4]{mhchem} % 用于绘制分子式
%    \end{macrocode}
%
% 用于生成可以被插入书签的LaTeX logo，TODO: 使用hologo创建|latex{}|命令
%    \begin{macrocode}
\RequirePackage{hologo} 
%    \end{macrocode}
%
% 
%    \begin{macrocode}
% Required to prevent page break right after a sectioning command
% \RequirePackage{needspace} 

% \RequirePackage{xspace} % Better print trailing whitespace
%    \end{macrocode}
%
% 针对编译引擎，使用不同的宏包重构下划线命令
%    \begin{macrocode}
\sys_if_engine_xetex:T
{
    \RequirePackage{microtype}

    % 加中文下划线，不能用于lualatex
    \RequirePackage{xeCJKfntef}
    \newcommand{\njuline}[1]{\CJKunderline{#1}}
}
\sys_if_engine_luatex:T{
    % 加中文下划线
    \RequirePackage{lua-ul}
    \newcommand{\njuline}[1]{\underLine{#1}\null}

    % emoji支持
    % \RequirePackage{emoji}
    % \setemojifont{Segoe~UI~Emoji} % windows
    % \setemojifont{Apple~Color~Emoji} % macos
    % \setemojifont{Noto~Color~Emoji}
    % For windows. 
    % Shipped with the best `grinning-face-with-sweat' support.
}
%    \end{macrocode}
%
% 生成用于测试的大段无意义英文文字
%    \begin{macrocode}
\RequirePackage{blindtext} 
\RequirePackage{zhlipsum} 
%    \end{macrocode}
%
% 加载以后展示内容边界
%    \begin{macrocode}
% \RequirePackage{showframe} 
%    \end{macrocode}
%
% 数学
%    \begin{macrocode}
\RequirePackage{amsmath} % Must be loaded before unicode-math
\RequirePackage{amsthm} % Mathematical environments
\RequirePackage{mathtools} % Mathematical tools to use with amsmath
\RequirePackage{thmtools} % Theorem styles
\RequirePackage[
    warnings-off={% 消除与mathtools合用产生的警告
        mathtools-colon,
        mathtools-overbracket}
        ]{unicode-math} % Math fonts in xetex or luatex
%    \end{macrocode}
%
% 配置图片路径
%    \begin{macrocode}
\RequirePackage{graphicx}
\DeclareGraphicsExtensions{.pdf,.eps,.jpg,.png}
\graphicspath{{figure/}} % 图片路径
\RequirePackage{wrapfig} % Wrap text around figures
%    \end{macrocode}
%
% 必须以该顺序加载以下包
%    \begin{macrocode}
\RequirePackage[hidelinks,bookmarksnumbered=true]{hyperref}
\RequirePackage[capitalise,nameinlink,noabbrev]{cleveref}
%    \end{macrocode}
%
% \subsection{字体设置}
%
% 加载设置字体所需的包\pkg{fontspec}
%    \begin{macrocode}
\RequirePackage{fontspec}
%    \end{macrocode}
%
% 
%
% \begin{variable}{\g_@@_latin_fontset_tl,\g_@@_cjk_fontset_tl}
% 用于存储所使用字体名称的全局变量
%    \begin{macrocode}
\tl_new:N \g_@@_latin_fontset_tl
\tl_new:N \g_@@_cjk_fontset_tl
%    \end{macrocode}
% \end{variable}
%
% 定义字体选项
%    \begin{macrocode}
\keys_define:nn { nju }
{
  customlatinfont   .choices:nn   =
  { gyre, macos, windows, none }
  { \tl_set_eq:NN \g_@@_latin_fontset_tl \l_keys_choice_tl },  
  customchinesefont .choices:nn   =
  { fandol, founder, macos, noto, windows, none }
  { \tl_set_eq:NN \g_@@_cjk_fontset_tl   \l_keys_choice_tl },
}
%    \end{macrocode}
%
% \begin{macro}{\ProcessKeysOptions}
% 在定义完全部设置以后从tex文件导言区输入参数
%    \begin{macrocode}
\ProcessKeysOptions { nju }
%    \end{macrocode}
% \end{macro}
%
% 如果未进行设置则使用系统默认字体覆盖输入参数
%
% \begin{variable}{\g_@@_load_system_fontset_bool}
% 用于判断是否需要载入系统预装字体的全局变量
%    \begin{macrocode}
\bool_new:N \g_@@_load_system_fontset_bool
%    \end{macrocode}
% \end{variable}
%
% 判断用户是否自定义了中英文字体，如果其中任意一种未被定义，
% 则使用系统预装字体覆盖字体选项
%    \begin{macrocode}
\tl_if_empty:NTF \g_@@_latin_fontset_tl
  { \bool_gset_true:N \g_@@_load_system_fontset_bool }  
{  
  \tl_if_empty:NT \g_@@_cjk_fontset_tl
    { \bool_gset_true:N \g_@@_load_system_fontset_bool }  
}
\bool_if:NT \g_@@_load_system_fontset_bool
{
  % 检测是否是 Windows
  \sys_if_platform_windows:TF
  {
    \tl_set:Nn \g_@@_latin_fontset_tl { windows }
    \tl_set:Nn \g_@@_cjk_fontset_tl   { windows }
  }
  {
    % 检测是否是 macOS
    \ctex_if_platform_macos:TF
    {
      \tl_set:Nn \g_@@_latin_fontset_tl { macos }
      \tl_set:Nn \g_@@_cjk_fontset_tl   { macos }
    }
    % 其余系统一律使用自由字体
    {
      \tl_set:Nn \g_@@_latin_fontset_tl { gyre }
      \tl_set:Nn \g_@@_cjk_fontset_tl { fandol }
    }
  }
}
%    \end{macrocode}
%
% 接下来逐个定义所需要使用的字库
%
% \begin{macro}{\@@_load_latin_font_windows:}
% Windows 西文字体
%    \begin{macrocode}
\cs_new_protected:Npn \@@_load_latin_font_windows:
{
  \setmainfont{Times~New~Roman}
  \setsansfont{Arial}
  \setmonofont{Courier~New}[Scale=MatchLowercase]
}
%    \end{macrocode}
% \end{macro}
%
%
% \begin{macro}{\@@_load_latin_font_macos:}
% macOS 西文字体
%    \begin{macrocode}
\cs_new_protected:Npn \@@_load_latin_font_macos:
{
\setmainfont{Times~New~Roman}
\setsansfont{Arial}
\setmonofont{Menlo}[Scale=MatchLowercase]
}
%    \end{macrocode}
% \end{macro}
%
%
% \begin{macro}{\@@_load_latin_font_gyre:}
% 开源的 gyre 西文字体
%    \begin{macrocode}
\cs_new_protected:Npn \@@_load_latin_font_gyre:
{
\setmainfont{texgyretermes}[
  Extension=.otf,
  UprightFont=*-regular,
  BoldFont=*-bold,
  ItalicFont=*-italic,
  BoldItalicFont=*-bolditalic]
\setsansfont{texgyreheros}[
  Extension=.otf,
  UprightFont=*-regular,
  BoldFont=*-bold,
  ItalicFont=*-italic,
  BoldItalicFont=*-bolditalic]
\setmonofont{texgyrecursor}[
  Extension=.otf,
  UprightFont=*-regular,
  BoldFont=*-bold,
  ItalicFont=*-italic,
  BoldItalicFont=*-bolditalic,
  Scale=MatchLowercase,
  Ligatures=CommonOff]
}
%    \end{macrocode}
% \end{macro}
%
% \begin{macro}{\@@_load_cjk_font_windows:}
% Windows 中文字体
%    \begin{macrocode}
\cs_new_protected:Npn \@@_load_cjk_font_windows:
{
  \setCJKmainfont{SimSun}[
    AutoFakeBold=2.17, 
    ItalicFont=KaiTi]
  \setCJKsansfont{SimHei}
  \setCJKmonofont{FangSong}
  \setCJKfamilyfont{zhsong}{SimSun}[AutoFakeBold=2.17]
  \setCJKfamilyfont{zhhei}{SimHei}
  \setCJKfamilyfont{zhfs}{FangSong}
  \setCJKfamilyfont{zhkai}{KaiTi}[AutoFakeBold=2.17]
  % \setCJKfamilyfont{zhnewhei}{Microsoft~YaHei}[
  %   BoldFont=Microsoft~YaHei~Bold]
}
%    \end{macrocode}
% \end{macro}
%
%
% \begin{macro}{\@@_load_cjk_font_macos:}
% macOS 字体
% 参考CTeX手册第9节 LuaLATEX 下的中文支持方式解决issue问题
%    \begin{macrocode}
\cs_new_protected:Npn \@@_load_cjk_font_macos:
{
  % 移除 does not contain script "CJK" 警告
  \msg_redirect_name:nnn {fontspec} {no-script} {info} 
  \setCJKmainfont{Songti~SC~Light}[
    BoldFont=Songti~SC~Bold,
    ItalicFont=Kaiti~SC,
    BoldItalicFont=Kaiti~SC~Bold]
  \setCJKsansfont{Heiti~SC~Light}[BoldFont=Heiti~SC~Medium]
  \setCJKmonofont{STFangsong}
  \setCJKfamilyfont{zhsong}{Songti~SC~Light}[BoldFont=Songti~SC~Bold]
  \setCJKfamilyfont{zhhei}{Heiti~SC~Light}[BoldFont=Heiti~SC~Medium]
  \setCJKfamilyfont{zhfs}{STFangsong}
  \setCJKfamilyfont{zhkai}{Kaiti~SC}[BoldFont=Kaiti~SC~Bold]
  \setCJKfamilyfont{zhnewhei}{PingFang~SC}
}
%    \end{macrocode}
% \end{macro}
%
%
% \begin{macro}{\@@_load_cjk_font_fandol:}
% Fandol 字体
%    \begin{macrocode}
\cs_new_protected:Npn \@@_load_cjk_font_fandol:
{
  % 移除 does not contain script "CJK" 警告
  \msg_redirect_name:nnn {fontspec} {no-script} {info} 
  \setCJKmainfont{FandolSong-Regular}[
    Extension=.otf,
    BoldFont=FandolSong-Bold,
    ItalicFont=FandolKai-Regular]
  \setCJKsansfont{FandolHei-Regular}[
    Extension=.otf,
    BoldFont=FandolHei-Bold]
  \setCJKmonofont{FandolFang-Regular}[Extension=.otf]
  \setCJKfamilyfont{zhsong}{FandolSong-Regular}[
    Extension=.otf,
    BoldFont=FandolSong-Bold]
  \setCJKfamilyfont{zhhei}{FandolHei-Regular}[
    Extension=.otf,
    BoldFont=FandolHei-Bold]
  \setCJKfamilyfont{zhfs}{FandolFang-Regular}[Extension=.otf]
  \setCJKfamilyfont{zhkai}{FandolKai-Regular}[
    Extension=.otf,
    AutoFakeBold=2.17]
}
%    \end{macrocode}
% \end{macro}
%
%
% \begin{macro}{\@@_load_cjk_font_founder:}
% 方正字库（简繁扩展）
% FZYouHK_508R \href{http://www.foundertype.com/index.php/FontInfo/index/id/244}{方正悠黑508R} 设计师非商免费，
% FZYouHK_511M \href{http://www.foundertype.com/index.php/FontInfo/index/id/244}{方正悠黑511M} 设计师非商免费
%    \begin{macrocode}
\cs_new_protected:Npn \@@_load_cjk_font_founder:
{
  \setCJKmainfont{FZSSK}[% 方正书宋
    Extension=.ttf,
    BoldFont=FZXBSK,% 方正小标宋
    ItalicFont=FZKTK]% 方正楷体
  \setCJKsansfont{FZXH1K}[% 方正细黑一
    Extension=.ttf,
    BoldFont=FZHTK]% FZHTK 方正黑体
  \setCJKmonofont{FZFSK}[Extension=.ttf]% 方正仿宋
  \setCJKfamilyfont{zhsong}
    {FZSSK}[
      Extension=.ttf,
      BoldFont=FZXBSK]
  \setCJKfamilyfont{zhhei}
    {FZHTK}[
      Extension=.ttf,
      AutoFakeBold=2.17]
  \setCJKfamilyfont{zhfs}
    {FZFSK}[Extension=.ttf]
  \setCJKfamilyfont{zhkai}
    {FZKTK}[Extension=.ttf]
  % \setCJKfamilyfont{zhnewhei}
  %   {FZYouHK_508R}[% 方正悠黑508R
  %     Extension=.ttf,
  %     BoldFont=FZYouHK_511M]% 方正悠黑511M
}
%    \end{macrocode}
% \end{macro}
%
%
% \begin{macro}{\@@_load_cjk_font_noto:}
% 思源字体
%    \begin{macrocode}
\cs_new_protected:Npn \@@_load_cjk_font_noto:
{
  % 思源宋体
  \setCJKmainfont[
    UprightFont=NotoSerifCJKsc-Regular,
    BoldFont=NotoSerifCJKsc-Bold,
    ItalicFont=NotoSerifCJKsc-Regular,
    BoldItalicFont=NotoSerifCJKsc-Bold,
    ItalicFeatures=FakeSlant,
    BoldItalicFeatures=FakeSlant]{Noto~Serif~CJK~SC}

  % 思源黑体
  \setCJKsansfont[
    UprightFont=NotoSansCJKsc-Regular,
    BoldFont=NotoSansCJKsc-Bold,
    ItalicFont=NotoSansCJKsc-Regular,
    BoldItalicFont=NotoSansCJKsc-Bold,
    ItalicFeatures=FakeSlant,
    BoldItalicFeatures=FakeSlant]{Noto~Sans~CJK~SC}

  % 包含于思源黑体
  \setCJKmonofont[
    UprightFont=NotoSansMonoCJKsc-Regular,
    BoldFont=NotoSansMonoCJKsc-Bold,
    ItalicFont=NotoSansMonoCJKsc-Regular,
    BoldItalicFont=NotoSansMonoCJKsc-Bold,
    ItalicFeatures=FakeSlant,
    BoldItalicFeatures=FakeSlant]{Noto~Sans~Mono~SC}

  \setCJKfamilyfont{zhsong}{Noto~Serif~CJK~SC}
  \setCJKfamilyfont{zhhei}{Noto~Sans~CJK~SC}
  \setCJKfamilyfont{zhfs}
    {FZFSK}[Extension=.ttf]
  \setCJKfamilyfont{zhkai}
    {FZKTK}[Extension=.ttf]
}
%    \end{macrocode}
% \end{macro}
%
% 载入设置的字体
%    \begin{macrocode}
\cs_new_protected:Npn \@@_load_font:
{
  \use:c { @@_load_latin_font_ \g_@@_latin_fontset_tl : }
  \use:c { @@_load_cjk_font_   \g_@@_cjk_fontset_tl   : }
  
  \NewDocumentCommand\songti{}{\CJKfamily{zhsong}}
  \NewDocumentCommand\heiti{}{\CJKfamily{zhhei}}
  \NewDocumentCommand\fangsong{}{\CJKfamily{zhfs}}
  \NewDocumentCommand\kaishu{}{\CJKfamily{zhkai}}
}
\@@_load_font:
%    \end{macrocode}
%
% 设置数学字体 (XITS, 或者 STIX, 与 Times New Roman 最为相近)
%    \begin{macrocode}
% \setmathfont{STIXTwoMath-Regular}[Extension = .otf]
\setmathfont{XITSMath-Regular}[
  BoldFont = XITSMath-Bold,
  Extension = .otf]
%    \end{macrocode}
%
%
% \subsection{页面样式}
%
% \subsubsection{页眉页脚}
%
% 
%    \begin{macrocode}
\RequirePackage{fancyhdr} % 调整页眉页脚
% 本科生页眉页脚 
\fancypagestyle{njuplain}{%
   \fancyhead{}               
   \fancyfoot[C]{\zihao{5}\thepage} % 页脚居中 五号新罗马体数字
}
%    \end{macrocode}
%
% 
%    \begin{macrocode}
% TODO: 研究生页眉页脚 
\fancypagestyle{njuheadings}{%
   \fancyhead{}               
   \fancyfoot[C]{\zihao{5}\thepage}        
}
%    \end{macrocode}
%
% 
%    \begin{macrocode}
\str_if_eq:NNTF {\l_@@_info_degree_tl} { ug } 
{
  % the header line
  \tl_set:Nn \headrulewidth {0pt}
  % the footer line
  \tl_set:Nn \footrulewidth {0pt}
  \AtBeginDocument{\pagestyle{njuplain}\flushbottom} % 本科无页眉页脚
}
{
  % the header line
  \tl_set:Nn \headrulewidth {1pt}
  % the footer line
  \tl_set:Nn \footrulewidth {0pt}
  \setlength{\headheight}{20pt}
  % \AtBeginDocument{\pagestyle{njuplain}} % 无页眉页脚
  \AtBeginDocument{\pagestyle{fancy}\flushbottom} % 研究生有页眉页脚
}
%    \end{macrocode}
%
% 章节标题格式
%    \begin{macrocode}
% 
\ctexset{
    chapter/format = \zihao{4}\heiti\centering\selectfont,
    chapter/beforeskip = 10pt,
    chapter/afterskip = 60pt,
    section/format = \zihao{4}\heiti\raggedright\selectfont,
    subsection/format = \zihao{4}\heiti\raggedright\selectfont,
    subsubsection/format = \zihao{4}\heiti\raggedright\selectfont,
}
%    \end{macrocode}
%
% 目录标题名称
%    \begin{macrocode}
\ctexset{
    contentsname = 目录,
    listfigurename = 插图清单, 
    listtablename = 表格清单,
}
%    \end{macrocode}
%
% \pkg{tocloft}定制目录文字格式
%    \begin{macrocode}
\RequirePackage[titles]{tocloft}
\cftsetpnumwidth{2em}
\renewcommand{\cftchapleader}{\cftdotfill{\cftchapdotsep}}
\renewcommand{\cftchapdotsep}{\cftdotsep}
\renewcommand{\cftchapfont}{\heiti\zihao{4}}
\setlength{\cftsecindent}{2em}
\setlength{\cftsubsecindent}{52pt}
\setlength{\cftsubsecnumwidth}{2em}
%    \end{macrocode}
%
% 
%
% 将目录插入书签
%    \begin{macrocode}
\renewcommand\tableofcontents{%
  \cleardoublepage
  \raggedbottom
  \begingroup
    \ctexset{
      contentsname = {目\hspace{2em}录},
      chapter/format = {\centering\songti\bf\zihao{3}\selectfont},
    }%
    \chapter*{\contentsname}%
  \endgroup
  \addcontentsline{toc}{chapter}{\contentsname}
  \vskip 20pt 
  \@starttoc{toc}%
}
\renewcommand\listoffigures{%
  \cleardoublepage
  \begingroup
  \ctexset{
    chapter/format = {\centering\songti\bf\zihao{3}\selectfont},
  }%
  \chapter*{\listfigurename}%
  \endgroup
  \addcontentsline{toc}{chapter}{\listfigurename}
  \vskip 20pt 
  \@starttoc{lof}%
}
\renewcommand\listoftables{%
  \cleardoublepage
  \begingroup
  \ctexset{
    chapter/format = {\centering\songti\bf\zihao{3}\selectfont},
  }%
  \chapter*{\listtablename}%
  \endgroup
  \addcontentsline{toc}{chapter}{\listtablename}
  \vskip 20pt 
  \@starttoc{lot}%
}
%    \end{macrocode}
%
% \begin{environment}{preface}
%
% 单独制作的前言页。
%    \begin{macrocode}
\NewDocumentEnvironment{preface}{}
{%
  \chapter*{前言}
  \addcontentsline{toc}{chapter}{前言}
}{\cleardoublepage}
%    \end{macrocode}
% \end{environment}
%
% \begin{environment}{acknowledgement}
% 单独制作的致谢页。
%    \begin{macrocode}
\NewDocumentEnvironment{acknowledgement}{}
{%
  \chapter*{致谢}
  \addcontentsline{toc}{chapter}{致谢}
}{\cleardoublepage}
%    \end{macrocode}
% \end{environment}
%
% 
% biblatex设置
%    \begin{macrocode}
\RequirePackage[
    style=gb7714-2015,
    %style=numeric-comp,
    %citestyle=authortitle-icomp,
    % citestyle=numeric-comp,
    %bibstyle=authoryear,
    % bibstyle=numeric,
    sorting=none,
    %sorting=nyt,
    %sortcites=true,
    %autocite=footnote,
    backend=biber, % Compile the bibliography with biber
    hyperref=true,
    backref=false,
    citecounter=true,
    pagetracker=true,
    citetracker=true,
    ibidtracker=context,
    autopunct=true,
    autocite=plain,
    % gbpub=false,         % Uncomment if you do NOT want '[S.l. : s.n.]' 
                           % in reference entries, GitHub Issue (#47)
    % gbnamefmt=lowercase, % Uncomment if you do NOT want uppercase author 
                           % names in reference entries, GitHub Issue (#23)
]{biblatex}
%    \end{macrocode}
%
% 
%    \begin{macrocode}
% Remove some unwanted entries from the bibliography
\AtEveryBibitem{
	\clearfield{abstract}
	\clearfield{issn}
	\clearfield{isbn}
	\clearfield{archivePrefix}
	\clearfield{arxivId}
	\clearfield{pmid}
	\clearfield{eprint}
	\ifentrytype{online}{}{\ifentrytype{misc}{}{\clearfield{url}}}
	% \ifentrytype{book}{\clearfield{doi}}{}
}
%    \end{macrocode}
%
% 
%    \begin{macrocode}
% 修改标签名称
\crefdefaultlabelformat{#2#1#3\,} % 默认在名称后面添加空格

\crefname{figure}{图}{图}
\crefname{table}{表}{表}
% \crefname{equation}{公式}{公式}
\crefformat{equation}{公式~#2#1#3~} % 删除公式编号的括号

\crefformat{chapter}{第#2#1#3章}
\crefformat{section}{第~#2#1#3~节}
\crefformat{subsection}{第~#2#1#3~小节}
\crefformat{subsubsection}{第~#2#1#3~小节}
% \crefname{chapter}{§}{§}
% \crefname{section}{§}{§}
% \crefname{subsection}{§}{§}
% \crefname{subsubsection}{§}{§}
\crefname{appendix}{附录}{附录}

% \crefname{definition}{定义}{定义}
% \crefname{axiom}{公理}{公理}
% \crefname{property}{性质}{性质}
% \crefname{proposition}{命题}{命题}
% \crefname{lemma}{引理}{引理}
% \crefname{corollary}{推论}{推论}
% \crefname{remark}{注解}{注解}
% \crefname{condition}{条件}{条件}
% \crefname{conclusion}{结论}{结论}
% \crefname{assumption}{假设}{假设}
%    \end{macrocode}
%
% soifjsojfosjfos
%    \begin{macrocode}
% 图表位置调整
\floatsetup[lstlisting]{ % Captions for lstlistings
	capposition=above,%
	margins=centering,%
	floatwidth=\textwidth%
}
\floatsetup[figure]{ % Captions for figures
	capposition=bottom,%
	margins=centering,%
	floatwidth=\textwidth%
}
\floatsetup[table]{ % Captions for tables
	capposition=above,%
	margins=centering,%
	floatwidth=\textwidth%
}
%    \end{macrocode}
%
% 
%    \begin{macrocode}
% 代码样式
\lstset{
	basicstyle=\ttfamily\linespread{1}\small\selectfont,
    keywordstyle=\bfseries,% use bold style for keywords
    commentstyle=\rmfamily\itshape,% use italic style for comments
    stringstyle=\ttfamily,% 字符串风格
    flexiblecolumns,% ?
    numbers=left,% left-aligned numbering
    showspaces=false,% hide markers for spaces
    showstringspaces=false,
    captionpos=t,% place the caption at the top
	% frame=lrtb,% show all four sides of the frame
	% linewidth=.8\textwidth,
	% breakatwhitespace=true,
	breaklines=true,
	xleftmargin=2em,xrightmargin=2em,% set the width of the code environment
}
%    \end{macrocode}
%
% 
%    \begin{macrocode}
\lstdefinestyle{LaTeX}{
  language=TeX,
  morekeywords={
    begin, caption, label, mathrm, frac, 
    toprule, midrule, bottomrule, includegraphics}
}
%    \end{macrocode}
%
% 
%    \begin{macrocode}
% 图表标题样式
\DeclareCaptionFont{songticap}{\zihao{5}\bf\songti}
\captionsetup{
    font=small,%
    labelfont=songticap,
	textfont=songticap,
	strut=no,%
	hypcap=true, % Links point to the top of the figure
	% indention=0pt, % Suppress indentation
	% % parindent=0pt, % Suppress space between paragraphs
	aboveskip=6pt, % Increase the space between the figure and the caption
	belowskip=6pt, % Increase the space between the caption and the table
}
%    \end{macrocode}
%
% 
%    \begin{macrocode}
% list configuration
\renewcommand{\labelitemi}{\tiny$\blacktriangleright$}
\renewcommand{\labelitemii}{\textbullet}

\setlist[itemize]{noitemsep}
\setlist[enumerate]{noitemsep}
\setlist[description]{noitemsep}
%    \end{macrocode}
%
% 
%    \begin{macrocode}
% 数学环境
\declaretheoremstyle[
    	%spaceabove=.5\thm@preskip,
    	%spacebelow=.5\thm@postskip,
    	headfont=\bf\songti,%\scshape,
    	notefont=\songti,% notebraces={ (}{)},
    	bodyfont=\songti,
    	%headformat={\NAME\space\NUMBER\space\NOTE},
    	headpunct={},
    	%postheadspace={.5em plus .1em minus .1em},
    	%prefoothook={\hfill\qedsymbol}
    ]{njuthm}

\theoremstyle{njuthm}
%    \end{macrocode}
%
% 
%    \begin{macrocode}
\let\oldproofname=\proofname
\renewcommand*{\proofname}{\rm\bf\songti{\oldproofname}} % 修改证明环境标题
%    \end{macrocode}
%
% 
%    \begin{macrocode}
\declaretheorem[
	name=算法,
	style=njuthm,
	refname={算法,算法},
	Refname={算法,算法},
	% numberwithin=section,
]{algorithm}
\declaretheorem[
	name=假设,
	style=njuthm,
	refname={假设,假设},
	Refname={假设,假设},
	% numberwithin=section,
]{assumption}
\declaretheorem[
	name=公理,
	style=njuthm,
	refname={公理,公理},
	Refname={公理,公理},
	% numberwithin=section,
]{axiom}
\declaretheorem[
	name=结论,
	style=njuthm,
	refname={结论,结论},
	Refname={结论,结论},
	% numberwithin=section,
]{conclusion}
\declaretheorem[
	name=条件,
	style=njuthm,
	refname={条件,条件},
	Refname={条件,条件},
	% numberwithin=section,
]{condition}
\declaretheorem[
	name=推论,
	style=njuthm,
	refname={推论,推论},
	Refname={推论,推论},
	% numberwithin=section,
]{corollary}
\declaretheorem[
	name=定义,
	style=njuthm,
	refname={定义,定义},
	Refname={定义,定义},
	% numberwithin=section,
]{definition}
\declaretheorem[
    	name=例,
    	style=njuthm,
    	refname={例,例},
    	Refname={例,例},
    	% numberwithin=section,
]{example}
\declaretheorem[
	name=引理,
	style=njuthm,
	refname={引理,引理},
	Refname={引理,引理},
	% numberwithin=section,
]{lemma}
\declaretheorem[
	name=性质,
	style=njuthm,
	refname={性质,性质},
	Refname={性质,性质},
	% numberwithin=section,
]{property}
\declaretheorem[
	name=命题,
	style=njuthm,
	refname={命题,命题},
	Refname={命题,命题},
	% numberwithin=section,
]{proposition}
\declaretheorem[
	name=注解,
	style=njuthm,
	refname={注解,注解},
	Refname={注解,注解},
	% numberwithin=section,
]{remark}
\declaretheorem[
	name=定理,
	style=njuthm,
	refname={定理,定理},
	Refname={定理,定理},
	numberwithin=section,
]{theorem}
%    \end{macrocode}
%
% \subsection{封面绘制}
%    \begin{macrocode}
% 拼合标题
\tl_const:Nn \l_@@_info_title_tl 
{
  \l_@@_info_title_a_tl
  \l_@@_info_title_b_tl
  \l_@@_info_title_c_tl
}
%    \end{macrocode}
% 双导师的处理
%% \begin{variable}{\l_@@_info_supv_full_tl,\l_@@_info_supv_full_en_tl}
% 用于存储导师姓名加职称的变量，旧版编译器不支持字符串中含有|\hspace{.5em}|这样的空白空间命令
%    \begin{macrocode}
\tl_new:N \l_@@_info_supv_full_tl
\tl_new:N \l_@@_info_supv_full_en_tl
%    \end{macrocode}
% \end{variable}
%
% 拼合导师
%    \begin{macrocode}
\tl_set:Nn \l_@@_info_supv_full_tl
{
  \l_@@_info_supv_a_tl\ 
  \l_@@_info_supv_a_title_tl\ \ 
  \l_@@_info_supv_b_tl\ 
  \l_@@_info_supv_b_title_tl
}
\tl_set:Nn \l_@@_info_supv_full_en_tl
{
  \l_@@_info_supv_a_en_tl\ 
  \l_@@_info_supv_a_title_en_tl\ \ 
  \l_@@_info_supv_b_en_tl\ 
  \l_@@_info_supv_b_title_en_tl
}
%    \end{macrocode}
% \begin{variable}{\c_@@_name_degree_tl,\c_@@_name_degree_en_tl}
% 用于存储学位名称的变量
%    \begin{macrocode}
\tl_new:N \c_@@_name_degree_tl
\tl_new:N \c_@@_name_degree_en_tl
%    \end{macrocode}
% \end{variable}
%
% 判断学位进行命令定义
%    \begin{macrocode}
\str_if_eq:NNTF {\l_@@_info_degree_tl} { ug } 
{
  % 本科
  \tl_const:Nn \l_@@_name_diploma_tl { 本\hfill 科\hfill }
  \tl_const:Nn \c_@@_name_title_tl { 题\hfill 目 }
  
  \tl_const:Nn \c_@@_cover_uline_len_a_tl { 250pt }
  \tl_const:Nn \c_@@_cover_uline_len_b_tl { 90pt }
  \tl_const:Nn \c_@@_cover_uline_font_tl { \songti }
  \tl_const:Nn \c_@@_cover_uline_style_tl { \bf }
  \tl_const:Nn \c_@@_cover_uline_bskip_tl {}

  \tl_const:Nn \c_@@_cover_box_len_tl { 4.2em }

  % 本科强制不打印国家图书馆封面
  \bool_set_false:N \g_@@_nlcover_bool
} 
{ 
  % 本科以外都是研究生 
  \tl_const:Nn \l_@@_name_diploma_tl { 研\hfill 究\hfill 生\hfill }
  \tl_const:Nn \c_@@_name_title_tl { 论\hfill 文\hfill 题\hfill 目 }

  \tl_const:Nn \c_@@_cover_uline_len_a_tl { 250pt }
  \tl_const:Nn \c_@@_cover_uline_len_b_tl { 14em }
  \tl_const:Nn \c_@@_cover_uline_font_tl { \kaishu }
  \tl_const:Nn \c_@@_cover_uline_style_tl {}
  \tl_const:Nn \c_@@_cover_uline_bskip_tl { \hspace{1em} }

  \tl_const:Nn \c_@@_cover_box_len_tl { 6em }

  % 研究生学位名称
  \str_if_eq:NNTF { \l_@@_info_degree_tl } { phd }
  {
    \tl_set:Nn \c_@@_name_degree_tl { 博士 }
    \tl_set:Nn \c_@@_name_degree_en_tl { Doctor~of~Philosophy }
  }
  {
    \tl_set:Nn \c_@@_name_degree_en_tl { Master }
    \str_if_eq:NNTF { \l_@@_info_degree_tl } { mg }
    { \tl_set:Nn \c_@@_name_degree_tl { 硕士 } }
    { \tl_set:Nn \c_@@_name_degree_tl { 硕士专业 } }
  }
}
%    \end{macrocode}
%
% \begin{macro}{\@@_cover_uline_a:n}
% 在封面绘制下划线
%    \begin{macrocode}
\cs_new:Npn \@@_cover_uline_a:n #1
{
  \uline{\makebox[\c_@@_cover_uline_len_a_tl]
    {\rm\c_@@_cover_uline_font_tl #1 }}
}
\cs_new:Npn \@@_cover_uline_b:n #1
{
  \uline{\makebox[\c_@@_cover_uline_len_b_tl]
    {\rm\c_@@_cover_uline_font_tl #1 }}
}
\cs_new:Npn \@@_cover_uline_nl:nn #1 #2
{
  \uline{\makebox[#1]
    {\rm\c_@@_cover_uline_font_tl #2 }}
}
%    \end{macrocode}
% \end{macro}
%
% 
% \begin{macro}{\@@_cover_box:n}
% 封面表格边框
%    \begin{macrocode}
\cs_new:Npn \@@_cover_box:n #1 
{
  \makebox[\c_@@_cover_box_len_tl][s]{
    #1\c_@@_cover_uline_bskip_tl}
}
%    \end{macrocode}
% \end{macro}
%
% 判断类型
% 
%    \begin{macrocode}
\str_if_eq:NNTF { \l_@@_info_type_tl } { thesis } 
{
  \tl_const:Nn \l_@@_info_type_tl_name 
    { \l_@@_name_diploma_tl 毕\hfill 业\hfill 论\hfill 文 }
}
{
  \tl_const:Nn \l_@@_info_type_tl_name 
    { \l_@@_name_diploma_tl 毕\hfill 业\hfill 设\hfill 计 }
}
%    \end{macrocode}
%
% \begin{macro}{\@@_cover_title_breakline:}
% 对多行标题进行自动判断，并输出换行后的结果
%    \begin{macrocode}
\cs_new_protected:Npn \@@_cover_title_breakline:
{
  \@@_cover_box:n {\c_@@_name_title_tl}
  & \@@_cover_uline_a:n 
  { \c_@@_cover_uline_style_tl \l_@@_info_title_a_tl } \\
  % 如果某行标题空则不输出接下来的若干行
  \tl_if_empty:NF \l_@@_info_title_b_tl
  {
    \tl_if_empty:NTF \l_@@_info_title_c_tl
      { 
        & \@@_cover_uline_a:n  
        {\c_@@_cover_uline_style_tl \l_@@_info_title_b_tl } \\ 
      }
      {
        & \@@_cover_uline_a:n  
        { \c_@@_cover_uline_style_tl \l_@@_info_title_b_tl } \\
        & \@@_cover_uline_a:n  
        { \c_@@_cover_uline_style_tl \l_@@_info_title_c_tl } \\
      }
  }
}
%    \end{macrocode}
% \end{macro}
%
% \begin{macro}{\@@_print_covertab:}
% 封面信息栏
%    \begin{macrocode}
\cs_new_protected:Npn \@@_print_covertab:
{%
  \str_if_eq:NNTF { \l_@@_info_degree_tl } { ug } 
  {
    \begin{tabular}{p{4.2em}c}
      \@@_cover_box:n {院\hfill 系}
      & \@@_cover_uline_a:n  {\l_@@_info_dept_tl} \\
      \@@_cover_box:n {专\hfill 业}
      & \@@_cover_uline_a:n  {\l_@@_major_tl} \\
      \@@_cover_title_breakline:
    \end{tabular}\\
  
    \begin{tabular}{p{4.2em}cp{4.2em}c}
      \@@_cover_box:n {年\hfill 级}
      & \@@_cover_uline_b:n {\l_@@_info_grade_tl}
      & \@@_cover_box:n {学\hfill 号}
      & \@@_cover_uline_b:n {\l_@@_info_id_tl}\\
    \end{tabular}\\
    
    \begin{tabular}{p{4.2em}c}
      \@@_cover_box:n {学\hfill 生\hfill 姓\hfill 名}
      & \@@_cover_uline_a:n  {\l_@@_info_author_tl}
    \end{tabular}\\
    \begin{tabular}{p{4.2em}cp{4.2em}c}
      \@@_cover_box:n {导\hfill 师}
      & \@@_cover_uline_b:n {\l_@@_info_supv_a_tl}
      & \@@_cover_box:n {职\hfill 称}
      & \@@_cover_uline_b:n {\l_@@_info_supv_a_title_tl}\\
  
      % 第二导师
      \tl_if_empty:NF \l_@@_info_supv_b_tl
      {
        \@@_cover_box:n {第\hfill 二\hfill 导\hfill 师}
        & \@@_cover_uline_b:n {\l_@@_info_supv_b_tl}
        & \@@_cover_box:n {职\hfill 称}
        & \@@_cover_uline_b:n {\l_@@_info_supv_b_title_tl}\\
      }
    \end{tabular}\\
  
    \begin{tabular}{p{4.2em}c}
      \@@_cover_box:n {提\hfill 交\hfill 日\hfill 期}
      & \@@_cover_uline_a:n  {\l_@@_submit_date_tl}\\
    \end{tabular}  
  } 
  { 
    \begin{tabular}{p{6em}c}    
      \@@_cover_title_breakline:
      \@@_cover_box:n {作\hfill 者\hfill 姓\hfill 名}
      & \@@_cover_uline_a:n  {\l_@@_info_author_tl}\\
      \@@_cover_box:n {专\hfill 业\hfill 名\hfill 称}
      & \@@_cover_uline_a:n  {\l_@@_major_tl}\\
      \@@_cover_box:n {研\hfill 究\hfill 方\hfill 向}
      & \@@_cover_uline_a:n  {\l_@@_field_tl}\\
      \@@_cover_box:n {指\hfill 导\hfill 教\hfill 师}
      & \@@_cover_uline_a:n  {\l_@@_info_supv_full_tl}\\
    \end{tabular}
  } 
}
%    \end{macrocode}
% \end{macro}
%
% 
%
% 
% \begin{macro}{\@@_print_cover_nl:}
% 按需绘制国家图书馆封面，修改自胡海星模板
% 
%    \begin{macrocode}
\cs_new_protected:Npn \@@_print_cover_nl:
{
  \thispagestyle{empty}
  \pdfbookmark[0]{国家图书馆封面}{nl}
  {
    % 顶端
    \hspace{0mm}
    \vskip -20mm \hskip -15mm
    \songti\zihao{-4}
    \makebox[40pt][l]{分类号}
    \@@_cover_uline_b:n {\l_@@_info_classif_tl}
    \hfill
    \makebox[40pt][l]{密级}
    \@@_cover_uline_b:n {\l_@@_info_seclv_tl}
    \vskip 10pt \hskip -15mm
    \makebox[40pt][l]{UDC}
    \@@_cover_uline_b:n {\l_@@_info_udc_tl}
  }

  % 中部
  \vskip\stretch{2}
  \begin{center}
    \def\ULthickness{1pt}
    {\kaishu\zihao{-0} 学\hspace{0.5em}位\hspace{0.5em}论\hspace{0.5em}文}
    {
      \kaishu\zihao{1}
      \vskip \stretch{1}
      \@@_cover_uline_b:n {\l_@@_info_title_a_tl}\\
      \@@_cover_uline_b:n {\l_@@_info_title_b_tl}\\
      \@@_cover_uline_b:n {\l_@@_info_title_c_tl}\\
    }
    \vskip \stretch{1}
    {\kaishu\zihao{4}（题名和副题名）}
    \vskip \stretch{1} \vskip 5mm
    {\kaishu\zihao{1}\uline{\makebox{\l_@@_info_author_tl}}}
    \vskip \stretch{1}
    {\kaishu\zihao{4}（作者姓名）}
  \end{center}

  % 底部
  \vskip\stretch{1}
  {
    \kaishu\zihao{4}
    \noindent 指导教师姓名、职务、职称、学位、单位名称及地址%
    \@@_cover_uline_nl:nn {94pt}{\l_@@_info_supv_a_tl}\par
    \noindent\@@_cover_uline_nl:nn {\textwidth}{%
    \l_@@_info_supv_cont_tl}\par
    \noindent 申请学位级别%
    \@@_cover_uline_nl:nn {9em}{\c_@@_name_degree_tl}%
    \noindent 专业名称%
    \uline{\hfill\l_@@_major_tl\hfill}\par% 需要调整下划线长度
    \noindent 论文提交日期%
    \@@_cover_uline_nl:nn {9em}{\l_@@_submit_date_tl}%
    论文答辩日期%
    \uline{\hfill\l_@@_defend_date_tl\hfill}\par% 需要调整下划线长度
    \noindent 学位授予单位和日期\uline{\hfill}\par
    \noindent\hfill 答辩委员会主席：%
    \@@_cover_uline_nl:nn {9em}{\l_@@_info_chairman_tl}\par
    \noindent\hfill 评阅人：%
    \@@_cover_uline_nl:nn {9em}{\l_@@_info_reviewer_a_tl}\par
    \noindent\hfill\@@_cover_uline_nl:nn {9em}{\l_@@_info_reviewer_b_tl}\par
    \noindent\hfill\@@_cover_uline_nl:nn {9em}{\l_@@_info_reviewer_c_tl}\par
    \noindent\hfill\@@_cover_uline_nl:nn {9em}{\l_@@_info_reviewer_d_tl}\par
    
    \begin{center}
      \kaishu\zihao{3}\hspace{2em} 年\hspace{1em} 月\hspace{1em} 日
    \end{center}
    \vskip -10mm
    }
  \cleardoublepage
}
%    \end{macrocode}
% \end{macro}
%
% 
% \begin{macro}{\@@_print_cover_ug:}
% 本科封面
%    \begin{macrocode}
\cs_new_protected:Npn \@@_print_cover_ug:
{
  \thispagestyle{empty}
  \pagenumbering{Roman}
  % Start
  \pdfbookmark{封面}{封面} % 将封面插入pdf书签
  \begin{spacing}{1.25}
  \vskip 0mm
  \hspace{-10mm}
  \includegraphics[height=3cm]{njulogo}\smallskip 
  \begin{center}
    \includegraphics[height=3.35cm]{njuname}
    \vskip 10mm 
    {\zihao{1}\makebox[9em][s]{\bf{\songti\l_@@_info_type_tl_name}}}
    \vfill
    \vskip\stretch{0}
    {\bgroup
    \kaishu\zihao{3}
    \def\tabcolsep{1pt}
    \def\arraystretch{1.5}
    % 绘制信息框
    \@@_print_covertab: 
    \egroup}
    \vfill
  \end{center}
  \end{spacing}
  \cleardoublepage
}
%    \end{macrocode}
% \end{macro}
%
% 
% \begin{macro}{\@@_print_cover_g:}
% 研究生封面
%    \begin{macrocode}
\cs_new_protected:Npn \@@_print_cover_g:
{
  \thispagestyle{empty}
  \pagenumbering{Roman}
  % Start
  \pdfbookmark{封面}{封面} % 将封面插入pdf书签
  \begin{spacing}{1.25}
  \begin{center}
    \hspace{0pt} \vskip 5mm
    \includegraphics[height=1.9cm]{njulogo}
    \vskip 10mm 
    \includegraphics[height=2cm]{njuname-large}
    \vskip 15mm 
    {\zihao{1}\makebox[10em][s]{\bf{\kaishu\l_@@_info_type_tl_name}}}
    \vskip 5mm
    {\zihao{1}\bf{\kaishu{（申请\c_@@_name_degree_tl 学位）}}}
    \par\vfill
    \vskip\stretch{0}
    {\bgroup
    \bf\kaishu\zihao{3}
    \def\tabcolsep{1pt}
    \def\arraystretch{1.5}
    \vskip 10mm
    % 绘制信息框
    \@@_print_covertab:
    \egroup}
    \vfill
    \vskip 10mm
    \bf\kaishu\zihao{4}\l_@@_submit_date_tl
    \vskip 15mm
  \end{center}
  \end{spacing}

  \newpage % 封面背面
  \thispagestyle{empty}
  \begin{spacing}{1.625}
    % TODO: 等待调整格式
    \hspace{0pt} \vfill
    {\bgroup
    \kaishu\zihao{3}
    \makebox[6em][s]{\bf\kaishu 学\hfill 号}：\MakeUppercase{\l_@@_info_id_tl}
    \par
    \makebox[6em][s]{\bf\kaishu 论文答辩日期}：\l_@@_defend_date_tl
    \par
    \makebox[6em][s]{\bf\kaishu 指\hfill 导\hfill 教\hfill 师}：\hspace{50mm}（签字）
    \par
    \egroup}
    \vskip 15mm
  \end{spacing}
  \cleardoublepage
}
%    \end{macrocode}
% \end{macro}
% 
% \begin{macro}{\@@_print_cover_en:}
% 研究生英文封面
%    \begin{macrocode}
\cs_new_protected:Npn \@@_print_cover_en:
{
  \thispagestyle{empty}
  \pdfbookmark{英文封面}{English~cover} % 将封面插入pdf书签
  \begin{center}
    \vspace*{20pt}
    \bf\sffamily\zihao{2}\l_@@_info_title_en_tl
    \vskip \stretch{1}
    \normalfont\rmfamily\zihao{4}{by}
    \vskip 3pt
    \bf\sffamily\zihao{4}\l_@@_info_author_en_tl
    \vskip\stretch{1}
    \normalfont\rmfamily\zihao{4}{Supervised~by}
    \vskip 3pt
    \normalfont\sffamily\zihao{4}\l_@@_info_supv_a_title_en_tl
    \hspace{.5em}\l_@@_info_supv_a_en_tl
    \vskip\stretch{1}
    \normalsize\rmfamily{%
      A~dissertation~submitted~to\\
      the~graduate~school~of~Nanjing~University\\
      in~partial~fulfilment~of~the~requirements~for~the~degree~of\\
      {\textsc{\c_@@_name_degree_en_tl}}\\
      in\\
      {\l_@@_major_en_tl}
    }
    \vskip\stretch{2}
    \includegraphics[width=2.5cm]{njulogo} \\
    \vskip 3mm
    \normalfont\l_@@_info_dept_en_tl\\
    {Nanjing~University}
    \vskip 30pt
    \normalfont\normalsize\l_@@_submit_date_en_tl
  \end{center}
  \normalfont
  \cleardoublepage
}
%    \end{macrocode}
% \end{macro}
%
% 
% 重定义maketitle生成封面
%    \begin{macrocode}
\tl_set:Nn \maketitle 
{%
  \str_if_eq:NNTF { \l_@@_info_degree_tl } { ug } 
  {
    \@@_print_cover_ug: % 本科封面
  } 
  { 
    \bool_if:NT \g_@@_nlcover_bool { \@@_print_cover_nl: } % 国家图书馆封面
    \@@_print_cover_g: % 研究生封面
    \@@_print_cover_en: % 英文封面
  } 
}
%    \end{macrocode}
%
% \subsection{摘要绘制}
%
% 中文关键词
%    \begin{macrocode}
\NewDocumentCommand \keywords {m} {%
  \par\vspace{2ex}\noindent%
  {\kaishu\zihao{-4}\makebox[4em][s]{关键词{：}}}~{#1}%
}
%    \end{macrocode}
%
% 
% 英文关键词
%
%    \begin{macrocode}
\NewDocumentCommand \englishkeywords {m} {%
  \par\vspace{2ex}\noindent%
  {KEYWORDS{:}}~~{#1}%
}
%    \end{macrocode}
%
%   
% 本科摘要环境
%
%    \begin{macrocode}
\cs_new_protected:Npn \@@_print_abstract_ug:
{
  % \pagestyle{plain}
  % \pagenumbering{Roman}
  % \phantomsection\addcontentsline{toc}{chapter}{中文摘要} % 将摘要插入目录和pdf书签
  \pdfbookmark[0]{中文摘要}{abstract-zh} % 将摘要插入pdf书签，与上一行不可共存
  \begin{center}
    \kaishu\zihao{-2}{\textbf{
      \uuline{南京大学本科生毕业论文（设计、作品）中文摘要}}}
  \end{center}
  {\bgroup
    \kaishu\zihao{-4}
    \tl_set:Nn \tabcolsep {0pt}
    \tl_set:Nn \arraystretch {0.8}
    \noindent
    题目： \l_@@_info_title_tl \\
    院系： \l_@@_info_dept_tl \\
    专业： \l_@@_major_tl \\
    本科生姓名： \l_@@_info_author_tl \\
    指导教师（姓名、职称）：\l_@@_info_supv_full_tl \\
    摘要：
    \egroup
  }
  \kaishu\zihao{-4}\par%
}
\cs_new_protected:Npn \@@_print_abstract_en_ug:
{
  \pagestyle{plain}
  % \phantomsection\addcontentsline{toc}{chapter}{英文摘要} % 将摘要插入目录和pdf书签
  \pdfbookmark[0]{英文摘要}{abstract-en} % 将摘要插入pdf书签，与上一行不可共存
  \begin{center}
      \kaishu\zihao{-2}{\textbf{\uuline{
        南京大学本科生毕业论文（设计、作品）英文摘要}}}
  \end{center}
  {
    \bgroup
    THESIS: ~~\l_@@_info_title_en_tl \\
    DEPARTMENT: ~~\l_@@_info_dept_en_tl \\
    SPECIALIZATION: ~~\l_@@_major_en_tl \\
    UNDERGRADUATE:~~\l_@@_info_author_en_tl \\
    MENTOR:~~\l_@@_info_supv_full_en_tl \\
    ABSTRACT:
    \egroup
  }
  \zihao{-4}\par%
}
%    \end{macrocode}
%
% 
% 研究生摘要环境
%    \begin{macrocode}
\cs_new_protected:Npn \@@_print_abstract_g:
{
  \pagestyle{plain}
  \pagenumbering{Roman}
  % \phantomsection\addcontentsline{toc}{chapter}{中文摘要} % 将摘要插入目录和pdf书签
  \pdfbookmark[0]{中文摘要}{abstract-zh} % 将摘要插入pdf书签，与上一行不可共存
  \begin{center}
    \kaishu\zihao{-2}{\textbf{\uuline{
      南京大学研究生毕业论文中文摘要首页用纸}}}
  \end{center}
  \bgroup
  \noindent
  \kaishu\zihao{4}
  \tl_set:Nn \tabcolsep {0pt}
  \tl_set:Nn \arraystretch {0.8}
  毕业论文题目：\hspace{0.5em}\njuline{\l_@@_info_title_tl\hfill}\\
  \njuline{\makebox[11em]{\l_@@_major_tl}}专业
  \njuline{\makebox[4em]{\l_@@_info_grade_tl}}级
  \str_if_eq:NNTF {\l_@@_info_degree_tl} { phd } {博}{硕}
  士生姓名：\njuline{\hfill\l_@@_info_author_tl\hfill}\\
  指导教师（姓名、职称）：\njuline{\hfill\l_@@_info_supv_full_tl\hfill}\par
  \egroup
  \kaishu\zihao{4}\par%
}
\cs_new_protected:Npn \@@_print_abstract_en_g:
{
  \pagestyle{plain}
  % \phantomsection\addcontentsline{toc}{chapter}{英文摘要} % 将摘要插入目录和pdf书签
  \pdfbookmark[0]{英文摘要}{abstract-en} % 将摘要插入pdf书签，与上一行不可共存
  \begin{center}
      \kaishu\zihao{-2}{\textbf{\uuline{
        南京大学研究生毕业论文英文摘要首页用纸}}}
  \end{center}
  {
    \bgroup
    \zihao{4}
    THESIS: ~~\l_@@_info_title_en_tl \\
    SPECIALIZATION: ~~\l_@@_major_en_tl \\
    POSTGRADUATE:~~\l_@@_info_author_en_tl \\
    MENTOR:~~\l_@@_info_supv_full_en_tl\par
    \egroup
  }
  \zihao{4}\par%
}
%    \end{macrocode}
%
% 
% 判断学位
%
%    \begin{macrocode}
\str_if_eq:NNTF {\l_@@_info_degree_tl} { ug } 
{   
  \NewDocumentEnvironment{abstract} {}
  {\@@_print_abstract_ug:}{\newpage}
  \NewDocumentEnvironment{englishabstract} {}
  {\@@_print_abstract_en_ug:}{\cleardoublepage}
}
{
  \NewDocumentEnvironment{abstract} {}
  {\@@_print_abstract_g:}{\newpage}
  \NewDocumentEnvironment{englishabstract} {}
  {\@@_print_abstract_en_g:}{\cleardoublepage}
}
%</class>
%    \end{macrocode}
%
%
%
%
% \end{implementation}
%
% \PrintIndex
